%% ma: L12-B17-0023
%% lop: 12
%% chuong: 5
%% bai: 17
%% dang: 12.17.15
%% loai: TN4
%% muc: NB
%% dap_an: "A"
%% nguon: "Ảnh giáo viên gửi 24/10/2026 (M06): Kỳ thi thử đợt IV, Thi thử TN THPT 2027 (Biên soạn: Aimbet), Phần 1 Câu 12"
%% kien_thuc: "Bài 17 (vị trí của điểm đối với mặt cầu), Bài 8 (khoảng cách giữa hai điểm)"
%% ghi_chu: "Đề gốc hỏi điểm nằm ngoài mặt cầu nhưng cả B, C, D đều nằm ngoài (IM² = 62, 61, 34 > 25); đã sửa câu hỏi thành 'nằm trong', đáp án A"
%% trang_thai: chinh-thuc
\begin{ex}
Trong không gian với hệ trục tọa độ $Oxyz$, cho mặt cầu $(S)$ có tâm $I(-2;3;5)$ và bán kính $R=5$. Điểm nào trong các điểm dưới đây nằm trong mặt cầu $(S)$?
\choice{\True $(-2;3;5)$}{$(5;0;3)$}{$(4;-1;2)$}{$(1;-1;8)$}
\loigiai{Bình phương khoảng cách đến tâm $I$: A: $0<25$; B: $49+9+4=62>25$; C: $36+16+9=61>25$; D: $9+16+9=34>25$. Chỉ điểm ở A nằm trong $(S)$.}
\end{ex}
